% !TeX root = ../main.tex
\section{Numerical implementation and verification}
\label{sec:implicit-ad-implementation}

The implementation evaluates the mineral volume and Biot coefficient at
integration points, updates plastic deformation when the yield condition is
reached, and assembles the phase balances.  Automatic differentiation (AD)
provides the derivatives of these calculations for Newton's method.

\subsection{Mineral state and plastic update}

At each integration point, the elastic or poroplastic constitutive update
solves the nonlinear scalar residual
\eqref{eq:verification-mineral-factors} using
\(x=\ln\bar J\) as the numerical unknown on the branch defined by
\eqref{eq:trace-mineral-stability-domain}.  A bracketing search locates the
root, and Newton iterations carried out with AD variables at that root supply
the first and second derivatives of the scalar equation.  After convergence,
the update recovers \(\bar J=\exp x\).  Intrinsic density follows from
\(\bar\rho_s=\bar\rho_{s0}/\bar J\), and the volume fraction is
\(\phi_s=\rho_s/\bar\rho_s\), using the current solid partial density.
The update then evaluates \eqref{eq:poroplastic-biot-correction} and checks
positive volumes, \(0<\phi_s<1\), positive mineral stiffness, and a positive
scalar derivative.

When plastic flow occurs, the seven equations
\eqref{eq:implicit-plastic-flow-increment}--\eqref{eq:discrete-consistency}
determine the six independent components of \(\mbf W\) and
\(\Delta\gamma\).  Each Newton iteration recomputes \(J^e\), mineral
volume, \(B\), and driving stress from the candidate plastic state.
For positive \(H\), the cohesion is evaluated at
\(\kappa^p_n+\Delta\gamma\); its derivative therefore enters the same
local consistency Jacobian.  Both \(\mbf F^p_n\) and \(\kappa^p_n\) remain
fixed during the current solve.
Forward-mode differentiation evaluates the local Jacobian with respect to
these seven unknowns.  Newton iteration first solves their values.  At the
converged root, the implicit-function theorem supplies their outer AD
derivatives by solving the same local Jacobian against the derivatives of
the residual with respect to the global fields, with previous-step history
fixed.  Scaling and squaring with a Taylor polynomial
through order 20 evaluates the matrix exponential using only algebraic
operations on AD scalars.  This construction avoids differentiating a spectral
decomposition when principal values coincide.

The elastic and poroplastic MOOSE materials share the scalar mineral solver
and closed-form Biot coefficient.  The mineral solve uses \(J^e=J/a^p\), while
the denominator of \eqref{eq:poroplastic-biot-correction} uses total \(J\).
The elastic specialization sets \(a^p=1\).
All current-state calculations retain their
AD derivatives.  The converged plastic factor is stored as ordinary values
for the next time step, while its current AD value remains in stress and
Biot coefficient calculations.  Differentiating the converged discrete update
supplies the derivative of the update that Newton iteration requires
\citep{simotaylor1985,miehe1996}.  MOOSE carries that derivative into the
assembled residual Jacobian \citep{lindsayetal2021ad}.

\subsection{Phase balances and finite-element residuals}
\label{sec:fluid-storage-and-transport}
\label{sec:weak-residuals-and-spaces}

The coupled boundary-value calculations use the elastic and poroplastic
specializations of the preceding constitutive model.  The solid mass balance determines partial
density, and the local mineral state determines solid volume fraction.
Fluid storage uses that pore volume and the barotropic EOS
\begin{align}
 \bar\rho_f(p)&=\bar\rho_{f0}
 \exp\left(\frac{p-p_0}{K_f}\right).
 \label{eq:barotropic-fluid-density}
\end{align}
Here \(K_f\) is the constant fluid bulk modulus and
\((p_0,\bar\rho_{f0})\) is its reference state.  Let \(\mbf w_f\)
be the fluid mass flux relative to the solid.  Its reference flux is
\(\mbf W_f=J\mbf F^{-1}\mbf w_f\).  The fluid balance is
\begin{equation}
 \left.\pder{}{t}\right|_{\mbf X}
 \left[J(1-\phi_s)\bar\rho_f\right]
 +\nabla_{\mbf X}\mathbin{\cdot}\mbf W_f=0.
 \label{eq:reference-fluid-mass-balance}
\end{equation}
For isotropic permeability \(\kappa\), dynamic viscosity \(\mu_f\), and
zero gravity, Darcy flow in reference coordinates is
\begin{equation}
 \mbf W_f=-\bar\rho_fJ\mbf F^{-1}
 \frac{\kappa}{\mu_f}\mbf F^{-T}\nabla_{\mbf X}p.
 \label{eq:reference-darcy-flux}
\end{equation}
The Biot coefficient enters momentum through the pressure term in stress.  The
same mineral volume determines fluid storage through \(\phi_s\).

Let \(\Omega_0\) be the solid reference domain, and let \(\delta\mbf u\),
\(\delta p\), and \(\delta\rho_s\) be test functions for displacement
\(\mbf u=\mbf x-\mbf X\), pressure, and solid partial density.  Let
\(\Gamma_t\) and \(\Gamma_f\) be the portions of the reference boundary with
prescribed nominal traction
\(\bar{\mbf t}_0=\mbf P\mbf N\) and prescribed outward referential fluid mass
flux \(\bar m_f=\mbf W_f\mathbin{\cdot}\mbf N\), respectively, where
\(\mbf N\) is the outward unit normal.  The displacement and pressure test
functions vanish on their corresponding essential boundaries.  With zero body
force, the residuals are
\begin{align}
 R_u(\delta\mbf u)
 &=\int_{\Omega_0}\nabla_{\mbf X}\delta\mbf u:
 \left(\mbf P^{\prime\prime}-BpJ\mbf F^{-T}\right)\dd V_0 \notag\\
 &\quad-\int_{\Gamma_t}\delta\mbf u\mathbin{\cdot}
 \bar{\mbf t}_0\dd A_0,
 \label{eq:momentum-weak-residual}\\
 R_p(\delta p)
 &=\int_{\Omega_0}\delta p\pder{}{t}
 \left[J(1-\phi_s)\bar\rho_f(p)\right]\dd V_0 \notag\\
 &\quad-\int_{\Omega_0}\nabla_{\mbf X}\delta p\mathbin{\cdot}
 \mbf W_f\dd V_0 \notag\\
 &\quad+\int_{\Gamma_f}\delta p\,\bar m_f\dd A_0,
 \label{eq:fluid-pressure-weak-residual}\\
 R_{\rho_s}(\delta\rho_s)
 &=\int_{\Omega_0}\delta\rho_s
 \left(\dot J\rho_s+J\dot\rho_s\right)\dd V_0.
 \label{eq:solid-mass-weak-residual}
\end{align}
The overdot follows the solid reference coordinate.  No boundary contribution
appears in \(R_{\rho_s}\) because the solid reference moves with the skeleton
and its mass balance contains no referential solid-flux divergence.  The
discrete fields use
\begin{equation}
 \mbf u^h\in[\mathcal Q_2]^2,
 \qquad p^h\in\mathcal Q_1,
 \qquad \rho_s^h\in\mathcal Q_2,
 \label{eq:finite-element-spaces}
\end{equation}
where \(\mathcal Q_k\) is the continuous tensor-product polynomial space.
The pressure uses continuous \(\mathcal Q_1\) elements without stabilization.  Solid
partial density is stored as \(\rho_s/\rho_{s0}\) for conditioning.

In the reported calculations, the MOOSE Backward Euler time integrator
supplies the rates of the displacement, pressure, and solid partial-density
fields.  The implementation then applies the chain rule directly to the solid
and fluid accumulations in \eqref{eq:solid-mass-weak-residual} and
\eqref{eq:fluid-pressure-weak-residual}, using rates derived from those field
rates.  Evaluating the fluid accumulation also requires \(\dot\phi_s\),
which depends on the mineral response and the plastic update.  The following
rates supply this constitutive contribution to
\eqref{eq:fluid-pressure-weak-residual}.

Since \(J^e=J/a^p\), differentiating the mineral EOS
\eqref{eq:verification-mineral-factors} along the current loading path gives
\begin{align}
 \left[\frac{K_s}{\bar J}
 +\left(1-\frac{K}{\phi_{s0}K_s}\right)p\right]\dot{\bar J}
 &=\frac{K}{\phi_{s0}}
 \left(\frac{\dot J}{J}-\frac{\dot a^p}{a^p}\right) \notag\\
 &\quad-\left(1-\frac{K}{\phi_{s0}K_s}\right)\bar J\dot p.
 \label{eq:poroplastic-mineral-volume-rate}
\end{align}
The identity \(\phi_s=\rho_s\bar J/\bar\rho_{s0}\) then combines this
mineral-volume rate with the solved solid partial-density rate to give
\begin{equation}
 \dot\phi_s=\phi_s\left(\frac{\dot\rho_s}{\rho_s}
 +\frac{\dot{\bar J}}{\bar J}\right).
 \label{eq:poroplastic-solid-fraction-rate}
\end{equation}
To evaluate the plastic contribution in
\eqref{eq:poroplastic-mineral-volume-rate}, the exponential update
\eqref{eq:plastic-exponential-update} supplies the logarithmic
distention rate over the same backward-Euler step,
\begin{equation}
 \frac{\ln a^p_n-\ln a^p_{n-1}}{\Delta t}
 =\frac{\beta\Delta\gamma}{\Delta t},
 \label{eq:discrete-plastic-distention-rate}
\end{equation}
which evaluates \(\dot a^p/a^p\) in
\eqref{eq:poroplastic-mineral-volume-rate}.  Here \(n\) denotes the current
time level and \(\Delta t\) the step.  The resulting solid-volume-fraction
rate enters the full water accumulation in
\eqref{eq:fluid-pressure-weak-residual}; its dependence on the converged
plastic increment remains in the outer AD Jacobian.  Positive plastic
distention increases pore volume at fixed total volume and pressure and
therefore contributes positively to water storage.

\subsection{Verification of the discrete residuals}

At a finite time step, the chain-rule rates used above are generally not equal
to backward differences of the complete nonlinear accumulations.  We therefore
measure the referential solid-mass drift and verify that it decreases under
time-step refinement.  A separate elastic BDF2 regression confirms that the same
chain-rule calculation accepts rates from a different MOOSE time integrator
without modification.

The Jacobian tests check the derivatives of these discrete residuals,
including the constitutive rates entering fluid storage.  PETSc forms a finite-difference approximation of
the residual Jacobian and compares it with the Jacobian assembled from the
MOOSE AD derivatives.  The relative and absolute Frobenius-norm differences
must remain below \(10^{-7}\) and \(10^{-5}\), respectively.  The elastic
coupled-field test includes momentum, solid mass, fluid storage, and Darcy
flow.  A separate active-plastic test applies the same comparison to algebraic
residuals containing the returned stress invariants, plastic distention, and
Biot coefficient at nonzero pressure and unequal principal stretches.  Each
finite-difference perturbation repeats the return mapping, so the comparison
tests the derivative of the converged local update carried into the outer
Newton Jacobian.  A further coupled poroplastic test includes all four field
residuals during compression and the subsequent hold.  A temporal predictor
moves the initial Newton guess away from the elastic--plastic switching
surface, and every assembled Jacobian is compared on the resulting path.
The finite-difference scale parameter is \(10^{-9}\) in the unit-scaled diagnostic;
the same relative and absolute tolerances apply.  At the switching surface
itself, perturbations can select different active branches, so a single smooth
Jacobian does not describe both sides.

An independent check perturbs total volume, pressure, solid partial density,
and plastic distention at sampled quadrature points, solves the scalar
mineral equation independently, and differences the complete water
accumulation.  It verifies the plastic storage contribution separately from
the assembled Jacobian.  Integrated Darcy outflow and the pressure-boundary
reaction are compared with water-mass change.  The reaction sums only nodes
carrying pressure degrees of freedom.  Time-step refinement measures the
mass drift introduced by the chain-rule time discretization.
